% Vietnamese translation for OI Public Library
% Source-PDF: IOI中国国家候选队论文/2002/孙方成.pdf
\documentclass[11pt]{article}
\usepackage{oipl-vn}

\title{Thuật toán và ứng dụng của đồ thị hai phía}
\author{Sun Fangcheng}
\date{}

\begin{document}
\maketitle

\origtitle{Algorithms and applications of bipartite graphs}
\sourcepath{IOI China candidate team papers / 2002 / Sun Fangcheng.pdf}

\section*{Tổng quan}

Bài viết này trước hết giới thiệu khái niệm ghép cặp, một quan hệ đặc biệt
trong đồ thị vô hướng, và định nghĩa đồ thị hai phía, một loại đồ thị đặc
biệt. Sau đó bài viết kết hợp hai khái niệm ấy để trình bày các thuật toán
hiệu quả cho ghép cặp cực đại theo số cạnh và ghép cặp tốt nhất trong đồ thị
hai phía. Đồng thời, thông qua quan hệ về số lượng giữa số ghép cặp cực đại
và số phủ nhỏ nhất của đồ thị hai phía, bài viết nêu rõ ưu thế riêng của đồ
thị hai phía so với đồ thị tổng quát.

\paragraph{Từ khóa.}
Đồ thị hai phía; ghép cặp; đường tăng; tập phủ; độ phức tạp thuật toán.

\tableofcontents

\section{Lời nói đầu}

Đồ thị hai phía là một loại đồ thị đặc biệt. Các đỉnh của nó luôn được chia
thành hai phần bù nhau; hai phần này thường được dùng để biểu diễn hai loại
sự vật khác nhau. Quan hệ cơ bản nhất giữa hai loại sự vật ấy chính là quan
hệ ghép cặp. Nếu biết kết hợp đồ thị hai phía và ghép cặp theo tình huống cụ
thể, ta có thể mở rộng đáng kể hướng suy nghĩ và tối ưu thuật toán. Nói tóm
lại, đồ thị hai phía, với tư cách một loại đồ thị đặc biệt, có ứng dụng rộng
rãi trong lập trình. Tính hiệu quả của nó giúp ta đưa ra lời giải hoàn chỉnh
cho một số trường hợp đặc biệt của các bài toán phức tạp.

\section{Khái niệm ghép cặp}

\paragraph{Định nghĩa 1.}
Cho đồ thị \(G=(V(G),E(G))\), và \(M\) là một tập con của \(E(G)\). Nếu hai
cạnh bất kỳ trong \(M\) đều không kề nhau, thì \(M\) được gọi là một
\term{ghép cặp} của \(G\). Hai đầu mút của một cạnh thuộc \(M\) được gọi là
được ghép đôi dưới \(M\).

Nếu một cạnh nào đó của ghép cặp \(M\) liên thuộc với đỉnh \(v\), ta nói
\(M\) bão hòa đỉnh \(v\), và \(v\) là đỉnh bão hòa đối với \(M\).

Giả sử \(M\) là một ghép cặp của đồ thị \(G\). Nếu trong \(G\) tồn tại một
đường đi đơn \(R\) mà các cạnh trên đường đi luân phiên là cạnh ghép thuộc
\(M\) và cạnh không ghép không thuộc \(M\), thì \(R\) được gọi là
\term{đường luân phiên}. Nếu hai đầu mút của \(R\) đều là đỉnh chưa bão hòa
đối với \(M\), thì đường luân phiên này được gọi là \term{đường tăng}.

Cho đồ thị \(G=(V(G),E(G))\), trong đó \(V(G)\) được chia thành hai tập đỉnh
con không rỗng và bù nhau \(X\) và \(Y\). Nếu một ghép cặp
\(M\subseteq E(G)\) của \(G\) bão hòa mọi đỉnh trong \(X\), nói cách khác
nếu toàn bộ đỉnh của \(X\) và một tập con đỉnh của \(Y\) xác định một quan
hệ một-một, thì \(M\) được gọi là một \term{ghép cặp hoàn bị} của \(G\).

Trong đa số trường hợp, nếu trực tiếp tìm ghép cặp cực đại hoặc ghép cặp tốt
nhất của một đồ thị từ góc nhìn đường tăng, hiệu quả thuật toán sẽ thấp. Vì
thế thuật toán ghép cặp cho đồ thị tổng quát còn phải liên quan đến định
nghĩa và xử lý ``nụ hoa''.\footnote{Trong ngữ cảnh này, ta gọi một chu trình
lẻ xuất hiện khi số cạnh ghép đạt cực đại là một nụ hoa.} Tuy nhiên, bài
viết này chủ yếu xét bài toán ghép cặp trên đồ thị hai phía, nên không khai
triển các nội dung ấy.

\section{Định nghĩa và nhận biết đồ thị hai phía}

\paragraph{Định nghĩa 2.}
Cho đồ thị \(G=(V,E)\). Nếu có thể chia \(V\) thành hai tập \(V_1\) và
\(V_2\), sao cho mỗi cạnh trong \(E\) có một đầu mút thuộc \(V_1\) và đầu
mút kia thuộc \(V_2\), thì đồ thị đó được gọi là \term{đồ thị hai phía},
cũng gọi là đồ thị hai phần hay đồ thị lưỡng phân.

Theo định nghĩa này, không có hai đỉnh nào trong cùng \(V_1\) hoặc cùng
\(V_2\) kề nhau. Đồ thị hai phía cũng có thể viết là
\(G=(V_1,V_2;E)\). Với tập đỉnh \(V\subseteq V_1\), ký hiệu \(P(V)\) là tập
các đỉnh trong \(V_2\) kề với ít nhất một đỉnh của \(V\).

\paragraph{Định nghĩa 3.}
Nếu trong đồ thị hai phía \(G\), mỗi đỉnh của tập con đỉnh bù \(V_1\) đều kề
với mọi đỉnh của \(V_2\), thì \(G\) được gọi là \term{đồ thị hai phía đầy
đủ}.

Với các đồ thị đơn giản, ta có thể nhận biết bằng mắt thường xem chúng có
phải đồ thị hai phía hay không. Nhưng với đồ thị phức tạp hơn, việc phán
đoán theo định nghĩa hoặc bằng cách vẽ lại trực quan rất bất tiện; vì vậy cần
một định lý nhận biết có hiệu quả.

\paragraph{Định lý 1.}
Một đồ thị vô hướng \(G\) là đồ thị hai phía khi và chỉ khi mọi chu trình của
\(G\) đều có độ dài chẵn. Nếu không có chu trình, có thể xem độ dài của mọi
chu trình là \(0\), và \(0\) được coi là số chẵn.

\paragraph{Chứng minh.}
Điều kiện cần: nếu \(G\) là đồ thị hai phía, thì mọi chu trình của \(G\) có
độ dài chẵn.

Lấy một chu trình tùy ý độ dài \(m\) trong \(G\):
\[
  C=v_{i_0},v_{i_1},v_{i_2},\ldots,v_{i_{m-1}},v_{i_0}.
\]
Vì \(G\) là đồ thị hai phía, ta có thể chia \(V\) thành hai tập bù nhau
\(V_1\) và \(V_2\). Do \(v_{i_0}\) và \(v_{i_1}\) kề nhau, giả sử
\[
  v_{i_0}\in V_1,\qquad v_{i_1}\in V_2.
\]
Tương tự suy ra
\[
  v_{i_2},v_{i_4},\ldots,v_{i_{m-2}}\in V_1,\qquad
  v_{i_3},v_{i_5},\ldots,v_{i_{m-1}}\in V_2.
\]
Nhìn vào chỉ số của các đỉnh, các đỉnh có chỉ số chẵn thuộc \(V_1\), còn các
đỉnh có chỉ số lẻ thuộc \(V_2\). Nay \(v_{i_{m-1}}\in V_2\), nên \(m-1\) là
số lẻ, tức \(m\) là số chẵn.

Điều kiện đủ: nếu mọi chu trình của \(G\) đều có độ dài chẵn, thì \(G\) phải
là đồ thị hai phía. Ta xét hai trường hợp.

\paragraph{Trường hợp 1: \(G\) liên thông.}
Chia tập đỉnh \(V\) của đồ thị thành hai tập con như sau:
\[
  V_1=\{v_i\mid \text{khoảng cách từ }v_i\text{ đến một đỉnh cố định }v_0
  \text{ là chẵn}\},\qquad
  V_2=V-V_1.
\]
Ta chứng minh \(V_1\) và \(V_2\) chính là hai tập đỉnh con bù nhau của \(V\).

Lấy tùy ý một cạnh \(e=(v_i,v_j)\) của \(G\). Nếu hai đầu mút \(v_i\) và
\(v_j\) của \(e\) đều nằm trong \(V_1\), thì cùng với các đường đi từ hai
đỉnh này đến \(v_0\), cạnh \(e\) tạo thành một chu trình
\[
  C=v_i,\ldots,v_0,\ldots,v_j,v_i.
\]
Do \(v_i,v_j\in V_1\), theo định nghĩa, khoảng cách từ \(v_i\) và \(v_j\) đến
\(v_0\) đều chẵn. Cộng thêm cạnh \(e\), độ dài chu trình \(C\) là lẻ, mâu
thuẫn với giả thiết. Vì thế \(v_i\) và \(v_j\) không thể cùng thuộc \(V_1\).

Nếu hai đầu mút của cạnh tùy ý \(e\) đều nằm trong \(V_2\), thì theo định
nghĩa, khoảng cách từ \(v_i\) và \(v_j\) đến \(v_0\) đều lẻ. Cộng thêm cạnh
\(e\), độ dài chu trình \(C\) vẫn là lẻ, cũng mâu thuẫn với giả thiết. Do đó
\(v_i\) và \(v_j\) cũng không thể cùng thuộc \(V_2\).

Vì vậy chỉ còn một khả năng duy nhất: hai đầu mút của \(e\), một đỉnh thuộc
\(V_1\) và đỉnh kia thuộc \(V_2\). Do \(e\) là cạnh tùy ý, theo định nghĩa
đồ thị hai phía, \(G\) là đồ thị hai phía.

\paragraph{Trường hợp 2: \(G\) không liên thông.}
Khi đó ta xét từng thành phần liên thông. Áp dụng chứng minh ở trên cho mỗi
thành phần liên thông của \(G\), rồi hợp các kết quả lại, ta được điều phải
chứng minh.

\section{Ghép cặp cực đại của đồ thị hai phía}

Theo định nghĩa, đồ thị hai phía vừa kế thừa các tính chất của đồ thị tổng
quát, vừa có thêm các tính chất riêng. Vì vậy thuật toán tìm ghép cặp cực đại
trên đồ thị hai phía đơn giản hơn thuật toán tương ứng trên đồ thị tổng quát,
và có phạm vi áp dụng rộng hơn. Có thể nói thuật toán ghép cặp cực đại của
đồ thị hai phía là nền tảng của mọi thuật toán trên đồ thị hai phía. Năm
1965, Edmonds đưa ra thuật toán Hungary để giải bài toán ghép cặp cực đại
trên đồ thị hai phía.

\begin{enumerate}
  \item Lấy một ghép cặp ban đầu \(M\) trong \(G\).
  \item Nếu mọi đỉnh trong \(X\) đều là đầu mút của cạnh thuộc \(M\), dừng:
  \(M\) là ghép cặp hoàn bị. Ngược lại, lấy một đỉnh \(u\in X\) không liên
  thuộc với cạnh nào trong \(M\), đặt \(S=\{u\}, T=\varnothing\).
  \item Nếu \(P(S)=T\), dừng: không có ghép cặp hoàn bị. Ngược lại, lấy
  \(y\in P(S)-T\).
  \item Nếu \(y\) là đầu mút của một cạnh thuộc \(M\), giả sử \(yz\in M\),
  đặt \(S\leftarrow S\cup\{z\}\), \(T\leftarrow T\cup\{y\}\), rồi quay lại
  bước 3. Ngược lại, lấy đường tăng \(R(u,y)\), đặt
  \(M\leftarrow M\oplus E(R)\), rồi quay lại bước 2.
\end{enumerate}

Phân tích thuật toán trên cho thấy độ phức tạp thời gian của thuật toán tìm
ghép cặp cực đại trong đồ thị hai phía là \(O(|E(G)|)\) cho mỗi quá trình
mở rộng theo cách mô tả này.

Trong đời sống thực tế, nhiều bài toán có thể chuyển thành bài toán ghép cặp
cực đại trên đồ thị hai phía, chẳng hạn bài toán phân công công việc. Ở đây
ta nêu một ví dụ tương đối mới: bài \textit{Chó con đi dạo}, lấy từ vòng
tuyển chọn Thượng Hải NOI 2001.

\begin{quote}
Grant thích dắt chú chó con Pandog đi dạo. Grant đi theo một tuyến đường cố
định với vận tốc không đổi; tuyến đường này có thể tự cắt. Pandog thích tham
quan các điểm cảnh vật dọc đường, nhưng sẽ gặp chủ tại \(N\) điểm đã cho.
Chó con và chủ xuất phát đồng thời từ điểm \((X_1,Y_1)\), và đồng thời hội
ngộ tại điểm \((X_n,Y_n)\). Vận tốc tối đa của chó con bằng hai lần vận tốc
của Grant. Khi chủ đi thẳng từ một điểm đến điểm kế tiếp, Pandog chạy đến
một điểm cảnh vật mà nó quan tâm. Trước mỗi lần gặp lại chủ, Pandog nhiều
nhất chỉ đi qua một điểm cảnh vật.

Nhiệm vụ là tìm cho Pandog một lộ trình, có thể trùng một phần với lộ trình
của chủ, sao cho nó tham quan được nhiều điểm cảnh vật nhất và vẫn kịp gặp
hoặc hội ngộ với chủ tại các vị trí đã cho.
\end{quote}

Nếu trực tiếp xét lộ trình của chó con, mỗi khi đến một điểm hội ngộ mới ta
đều phải xét phương án di chuyển của chó. Tổng lượng tính toán là
\(P_m^n\); thuật toán có độ phức tạp mũ như vậy rõ ràng kém hiệu quả và
không thực tế. Đồng thời, mô hình của đề bài có liên hệ với mô hình hoán vị,
và tính hậu hiệu quá cao về căn bản phủ định khả năng dùng quy hoạch động.
Ngay cả khi dùng mô hình luồng mạng, nếu chỉ trực tiếp xây dựng từ góc nhìn
của chó thì cũng rất khó cài đặt. Nhưng nếu nắm được một số mô tả then chốt
của đề và trừu tượng hóa chúng thành mô hình đồ thị hai phía, ta sẽ có hiệu
quả tốt hơn.

Vì ``khi chủ đi thẳng từ một điểm đến điểm kế tiếp, Pandog chạy đến một điểm
cảnh vật mà nó quan tâm'', và lộ trình di chuyển của chủ đã cố định, việc tìm
một lộ trình có thể xem là quyết định lựa chọn của chó trên từng đoạn đường:
chạy đến điểm cảnh vật nào, hoặc đi trùng với lộ trình của chủ. Rõ ràng, nếu
trừu tượng hóa mỗi đoạn trong lộ trình của chủ thành một đỉnh, thì tính chất
của các đỉnh này khác với tính chất của các điểm cảnh vật mà chó muốn đến.
Vì vậy các đỉnh trong đồ thị tự nhiên được chia thành hai loại, và điều đề
bài yêu cầu chính là quan hệ giữa hai loại đỉnh ấy.

Trên một đoạn thẳng, chó chỉ có thể chạy đến một điểm cảnh vật, nên việc một
điểm cảnh vật có đến được hay không chỉ liên quan đến khoảng cách từ điểm ấy
đến hai đầu mút của đoạn thẳng. Một đoạn thẳng chỉ có thể liên hệ với một
điểm cảnh vật, và một điểm cảnh vật chỉ có ý nghĩa ở lần đầu tiên được thăm,
nên bất kỳ hai quan hệ liên hệ nào cũng không giao nhau về đối tượng sử dụng.
Điều này rõ ràng thỏa định nghĩa ghép cặp. Vì thế, điều cần tìm của đề chính
là ghép cặp cực đại giữa hai tập đỉnh trong một đồ thị gồm hai loại đỉnh.

Quá trình xây dựng đồ thị hai phía như sau. Giả sử có \(n\) điểm định vị và
\(m\) điểm cảnh vật:
\[
  V_1=\{\overrightarrow{P_iP_{i+1}}\mid 1\le i<n\},\qquad
  V_2=\{Q_i\mid 1\le i\le m\},\qquad
  V(G)=V_1\cup V_2.
\]
Với đoạn có hướng \(\overrightarrow{P_iP_{i+1}}\) và điểm cảnh vật \(Q_k\),
khi và chỉ khi
\[
  D(P_i,Q_k)+D(Q_k,P_{i+1})\le 2D(P_i,P_{i+1}),
\]
thì cạnh \((\overrightarrow{P_iP_{i+1}},Q_k)\) thuộc \(E(G)\). Ở đây
\[
  D(A,B)=\sqrt{(X_A-X_B)^2+(Y_A-Y_B)^2}.
\]
Đối chiếu với định nghĩa đồ thị hai phía, đây là một đồ thị hai phía điển
hình. Sau đó ta dùng thuật toán Hungary tìm ghép cặp cực đại trên đồ thị hai
phía để thu được số ghép cặp cực đại và phương án ghép cặp, cuối cùng sinh
dữ liệu ra theo yêu cầu đề bài.

\section{Bài toán tập phủ nhỏ nhất của đồ thị hai phía}

Thuật toán ghép cặp cực đại của đồ thị hai phía có ứng dụng rộng rãi trong
giải bài toán thực tế. Trường hợp điển hình nhất là dùng nó để giải bài toán
tập phủ nhỏ nhất của đồ thị hai phía. Chẳng hạn bài \textit{Knights} của BOI
2001 có mô tả như sau:

\begin{quote}
Cho một bàn cờ quốc tế kích thước \(n\times n\), trong đó một số ô đã bị lấy
đi. Nhiệm vụ của bạn là xác định cách đặt nhiều quân mã nhất trên bàn cờ sao
cho chúng không tấn công nhau. Kích thước bàn cờ thỏa \(n\le 200\).
\end{quote}

Một quân mã đặt tại ô \(S\) tấn công tám ô có độ lệch tọa độ
\((\pm 1,\pm 2)\) và \((\pm 2,\pm 1)\), nếu các ô ấy còn nằm trên bàn cờ.

Đề yêu cầu đặt nhiều quân mã nhất và không để chúng tấn công nhau. Các bài
toán có mô tả tương tự rất phổ biến; quen thuộc nhất là bài toán tám hậu.
Bài tám hậu là ví dụ kinh điển của thuật toán quay lui. Vì ở mỗi hàng gần
như mọi ô đều phải được thử một lần, độ phức tạp thời gian là \(O(n!)\), và
đây còn là một bài toán NP kinh điển. Nếu giải bài này theo tư tưởng của tám
hậu, thách thức sẽ rất nghiêm trọng.

Trước hết, cách mở rộng trạng thái của tám hậu đơn giản hơn bài này rất
nhiều. Vì phạm vi tấn công của hậu là cả một hàng, một cột hoặc một đường
chéo, khi mở rộng ta có thể tránh rất nhiều trạng thái vô hiệu. Trong bài
này, phạm vi tấn công tương đối phức tạp, nên việc trực tiếp tránh trạng thái
vô hiệu trong quá trình tìm kiếm khó hơn, và hiệu quả thuật toán tự nhiên
không lý tưởng. Thứ hai, quy mô bài toán đạt tới \(200\), đã vượt rất xa
phạm vi có thể giải của tám hậu. Đồng thời, do phạm vi khống chế của quân mã
nhỏ, số phương án khả thi chắc chắn tăng lên; còn quay lui hay các thuật toán
tìm kiếm khác đều dựa trên việc tìm phương án để thu được lời giải. Tư tưởng
cơ bản vẫn là vét cạn, nên số lời giải tăng sẽ làm hiệu quả giảm mạnh. Trong
khi đó mục tiêu của bài này lại là số quân mã đặt được nhiều nhất; thuật toán
tìm kiếm rõ ràng không đủ đúng trọng tâm.

Tiếp theo xét bài \textit{Pháo binh trận địa} của NOI 2001. Trong bài ấy,
người ta dùng cách tách trạng thái và cuối cùng thực hiện được quy hoạch
động. Nhưng trong quá trình quy hoạch động, nếu xét theo chiều rộng, độ phức
tạp vẫn là mũ. So với \textit{Pháo binh trận địa}, sự không đều của phạm vi
tấn công cũng là bất lợi của bài này, khiến quá trình chuyển trạng thái khó
hơn nhiều. Quan trọng hơn, mấu chốt để quy hoạch động giải được
\textit{Pháo binh trận địa} nằm ở sự bất đối xứng rất mạnh giữa chiều dài và
chiều rộng. Trong bài này, bàn cờ là hình vuông đúng nghĩa, nên tư tưởng
tách trạng thái rất khó thực hiện, cả về thời gian lẫn không gian. Do đó quy
hoạch động cũng không phù hợp với bài này. Nói cách khác, khi suy luận bằng
những hướng thông thường, ta không thu được phương pháp khả thi cho bài toán.

Bây giờ phân tích mô hình của bài từ góc độ xây dựng đồ thị. Xem mỗi ô của
bàn cờ là một đỉnh, và xem quan hệ hai ô có thể đi trực tiếp tới nhau, tức
có thể tấn công nhau, là một cạnh. Điều đề bài yêu cầu là tập độc lập lớn
nhất trong đồ thị ấy. Vì với bất kỳ đồ thị nào, tập độc lập lớn nhất và tập
phủ nhỏ nhất là hai tập bù nhau, bài này cũng là một bài toán đồ thị tìm tập
phủ nhỏ nhất. Như đã biết, bài toán tập phủ nhỏ nhất trên đồ thị tổng quát là
một bài toán NPC. Vậy bài này có thoát được tình thế ấy không? Điều đó chỉ
có thể phân tích từ tính đặc thù của bài.

Kiến thức thông thường cho ta biết trên bàn cờ quốc tế, ô trắng và ô đen xen
kẽ nhau; một quân mã chỉ có thể nhảy từ ô trắng sang ô đen, hoặc từ ô đen
sang ô trắng. Vì tính đặc biệt này, phạm vi tấn công không đều của quân mã
lại đem đến thuận lợi cho việc giải bài. Một quân mã ở một vị trí chỉ có thể
tấn công các ô khác màu với vị trí của nó; đây là điều kiện cần. Nói cách
khác, giữa hai ô bất kỳ cùng màu không tồn tại cạnh. Vì vậy bàn cờ có thể
được chia thành hai phần theo màu, và trong cùng một phần, hai đỉnh bất kỳ
đều không thể tấn công nhau. Điều đó có nghĩa bàn cờ mà đề mô tả, khi định
nghĩa trên quan hệ tấn công, là một đồ thị hai phía.

Đồng thời, nhờ tính chất riêng của đồ thị hai phía, bài toán tìm số phủ nhỏ
nhất của nó có thuật toán đa thức. Giữa số ghép cặp cực đại và số phủ nhỏ
nhất của đồ thị hai phía tồn tại một quan hệ xác định. Năm 1931, König đưa
ra định lý về quan hệ giữa ghép cặp cực đại và phủ nhỏ nhất như sau.

\paragraph{Định lý 2.}
Giả sử \(r\) là số ghép cặp cực đại của đồ thị hai phía \(G\), và \(s\) là
số phủ nhỏ nhất của ma trận kề của nó. Khi đó \(r=s\).

Để chứng minh định lý này, ta dùng định lý hôn nhân Hall: giả sử \(G\) là
một đồ thị hai phía có hai tập đỉnh lần lượt là \(V\) và \(W\). Khi đó tồn
tại một ghép cặp hoàn bị của \(G\) khi và chỉ khi
\[
  |S|\le |P(S)|\qquad\text{với mọi }S\subseteq V.
\]
Chứng minh định lý Hall được cho trong phụ lục. Sau đây là chứng minh định
lý ghép cặp cực đại bằng phủ nhỏ nhất.

Vì mỗi phần tử khác \(0\) không cùng hàng, không cùng cột với nhau cần một
hàng hoặc một cột để phủ, nên \(r\) phần tử khác \(0\) đôi một không cùng
hàng, không cùng cột cần \(r\) hàng hoặc cột để phủ. Trong khi đó \(s\) hàng
và cột khác nhau đã phủ toàn bộ phần tử khác \(0\) của ma trận \(A\), nên
tự nhiên cũng phủ \(r\) phần tử khác \(0\) nói trên. Do đó \(s\ge r\). Ta
tiếp tục chứng minh \(r\ge s\).

Không mất tính tổng quát, giả sử phủ nhỏ nhất phủ \(c\) hàng và \(d\) cột
của \(A\), tức \(s=c+d\). Gọi tập đỉnh tương ứng với \(c\) hàng ấy là
\(X_c\), phần còn lại là \(X-X_c\); gọi tập đỉnh tương ứng với \(d\) cột ấy
là \(Y_d\), phần còn lại là \(Y-Y_d\). Sắp xếp lại các hàng và cột của ma
trận \(A\) như sơ đồ sau:
\[
\begin{array}{c|cc}
 & Y-Y_d & Y_d\\ \hline
X_c & A_{11} & A_{12}\\
X-X_c & A_{21} & A_{22}
\end{array}
\]
Rõ ràng mọi hàng của \(A_{11}\) đều được phủ, mọi cột của \(A_{22}\) đều
được phủ, mọi phần tử của \(A_{12}\) được phủ bởi cả hàng lẫn cột, còn
\(A_{21}=0\).

Ta chứng minh tồn tại ghép cặp hoàn bị từ \(X_c\) đến \(Y-Y_d\). Trong
\(A_{11}\), lấy tùy ý \(V'\) hàng; các phần tử khác \(0\) trong những hàng
này ít nhất phải phân bố trên \(|V'|\) cột khác nhau. Nếu không, ta có thể
không phủ \(|V'|\) hàng này mà phủ số cột ít hơn ấy, và toàn bộ phần tử khác
\(0\) trong \(A\) vẫn được phủ. Khi đó số phần dùng để phủ nhỏ hơn ban đầu,
mâu thuẫn với giả thiết phủ ban đầu là nhỏ nhất. Nói cách khác, với mọi tập
con \(V'\) của \(X_c\), luôn có \(|P(V')|\ge |V'|\). Theo định lý hôn nhân
Hall, tồn tại ghép cặp hoàn bị \(M_1\) từ \(X_c\) đến \(Y-Y_d\), và
\(|M_1|=c\). Tương tự, cũng tồn tại ghép cặp hoàn bị \(M_2\) từ \(Y_d\) đến
\(X-X_c\), và \(|M_2|=d\). Hợp \(M_1\cup M_2\) vẫn là một ghép cặp của
\(G\), đồng thời
\[
  |M_1\cup M_2|=c+d=s.
\]
Vì vậy số ghép cặp cực đại \(r\ge s\). Chứng minh hoàn tất.

Nhờ định lý này, các thuật toán liên quan đến tập phủ nhỏ nhất, tập độc lập
lớn nhất, v.v. của đồ thị hai phía đều có thể thu được gián tiếp thông qua
việc tìm số ghép cặp cực đại của đồ thị hai phía.

\section{Bài toán ghép cặp tốt nhất của đồ thị hai phía}

Phần trước đã đưa ra thuật toán tìm ghép cặp cực đại của đồ thị hai phía.
Nhưng trong công việc thực tế, hiệu suất thường là yếu tố không thể thiếu
khi xét bài toán. Chẳng hạn bài toán phân công công việc với hiệu suất làm
việc khác nhau, hoặc bài \textit{Chó con đi dạo} sau khi thêm trọng số cho
các điểm cảnh vật, không thể chỉ dùng ghép cặp cực đại để giải. Dưới đây là
một bài toán phân công công việc có trọng số.

\begin{quote}
Một công ty có các nhân viên \(x_1,x_2,\ldots,x_n\), họ cần làm các công
việc \(y_1,y_2,\ldots,y_n\). Mỗi người có thể làm một vài công việc trong số
đó, và với mỗi công việc có một hiệu suất cố định. Hỏi có thể tìm một phương
án phân công thích hợp sao cho tổng hiệu suất là cao nhất hay không. Yêu cầu
mỗi người chỉ tham gia một công việc, đồng thời mỗi công việc cũng phải do
một người độc lập hoàn thành. Không yêu cầu mọi người đều có công việc.
\end{quote}

Với dữ liệu trong bảng sau, nếu một nhân viên hoàn toàn không thể hoàn thành
một công việc thì hiệu suất được ghi là \(0\):
\[
\begin{array}{c|cccc}
 & \text{A} & \text{B} & \text{C} & \text{D}\\ \hline
\text{Giáp} & 1 & 10 & 0 & 0\\
\text{Ất} & 0 & 1 & 0 & 0\\
\text{Bính} & 0 & 0 & 1 & 0\\
\text{Đinh} & 0 & 0 & 0 & 1
\end{array}
\]
Nếu dùng thuật toán ghép cặp có số cạnh lớn nhất, ta có thể nhận được lời
giải: Giáp--A, Ất--B, Bính--C, Đinh--D, với tổng giá trị cuối cùng là \(4\).
Nhưng nếu không ghép dự án A với Giáp mà dùng Ất--A, Bính--C, Đinh--D, tuy
số cạnh ghép ít hơn nhưng tổng giá trị cuối cùng đạt \(12\). Thực tế chứng
tỏ thuật toán ghép cặp cực đại không có tác dụng thực chất với các bài toán
đặt hiệu ích lên hàng đầu như vậy.

Vậy có con đường hữu hiệu nào khác để giải loại bài toán này không? Vì mỗi
người chỉ có thể làm một công việc, và mỗi công việc cũng chỉ có thể do một
người làm, tính hậu hiệu trong quyết định phương án là rất rõ. Nếu muốn thoát
khỏi rắc rối hậu hiệu bằng cách thay đổi mô tả trạng thái, ta sẽ gặp vấn đề
chí mạng của quy hoạch động là thảm họa số chiều. Do đó quy hoạch động không
thể dùng để giải loại bài toán này.

Một phương pháp khả thi là nhánh cận. Tức là thông qua lựa chọn hợp lý hàm
đánh giá, trên cơ sở tìm kiếm theo chiều sâu, liên tục cập nhật kết quả tối
ưu để giảm tìm kiếm vô hiệu và nâng cao hiệu quả. Tuy nhiên, mặc dù nhánh
cận có hiệu quả khá tốt trong ứng dụng quy mô vừa và nhỏ, bản thân nó vẫn có
độ phức tạp mũ, khiến hiệu quả giảm nhanh khi quy mô tăng.

Thực ra, loại bài toán này chính là bài toán ghép cặp tốt nhất trong lý
thuyết đồ thị. Mô tả như sau:
\[
  G=(V_1,V_2;E)
\]
là đồ thị hai phía đầy đủ có trọng số,
\[
  V_1=\{x_1,x_2,\ldots,x_n\},\qquad
  V_2=\{y_1,y_2,\ldots,y_n\},
\]
với trọng số \(w(x_iy_j)\ge 0\). Cần tìm ghép cặp hoàn bị có tổng trọng số
lớn nhất; ghép cặp như vậy được gọi là \term{ghép cặp tốt nhất}. Vì định
nghĩa yêu cầu đồ thị hai phía đầy đủ có trọng số, xét từ góc độ thuật toán,
mục đích chính là bảo đảm chắc chắn tồn tại ghép cặp hoàn bị. Vì vậy, khi
một số cạnh trong đồ thị ban đầu không tồn tại, ta thường gán cho chúng trọng
số \(0\), \(1\), hoặc một giá trị khác không ảnh hưởng đến kết quả cuối cùng,
để làm cho đồ thị trở thành đầy đủ. Sau khi chuẩn bị như vậy, trước khi giải
cụ thể loại bài toán này, ta cần biết một số khái niệm liên quan.

\paragraph{Định nghĩa 4.}
Một ánh xạ \(l:V(G)\to R\) thỏa mãn
\[
  \forall x\in V_1,\quad \forall y\in V_2,\quad
  l(x)+l(y)\ge w(xy)
\]
được gọi là \term{nhãn đỉnh khả thi} của đồ thị hai phía \(G\). Đặt
\[
  E_l=\{xy\mid xy\in E(G),\ l(x)+l(y)=w(xy)\}.
\]
Đồ thị con sinh của \(G\) lấy \(E_l\) làm tập cạnh được gọi là
\term{đồ thị bằng nhau}, ký hiệu \(G_l\).

Nhãn đỉnh khả thi luôn tồn tại, chẳng hạn
\[
l(v)=
\begin{cases}
\max_{y\in V_2} w(xy), & v=x\in V_1,\\
0, & v=y\in V_2.
\end{cases}
\]

\paragraph{Định lý 3.}
Ghép cặp hoàn bị của \(G_l\) chính là ghép cặp tốt nhất của \(G\).

\paragraph{Chứng minh.}
Giả sử \(M^*\) là một ghép cặp hoàn bị của \(G_l\). Vì \(G_l\) là đồ thị con
sinh của \(G\), nên \(M^*\) cũng là ghép cặp hoàn bị của \(G\). Tập các đầu
mút của các cạnh trong \(M^*\) chứa mỗi đỉnh của \(G\) đúng một lần, do đó
\[
  W(M^*)=\sum_{e\in M^*}w(e)=\sum_{v\in V(G)}l(v).
\]
Mặt khác, nếu \(M\) là một ghép cặp hoàn bị bất kỳ trong \(G\), thì
\[
  W(M)=\sum_{e\in M}w(e)\le \sum_{v\in V(G)}l(v).
\]
Vì vậy \(W(M^*)\ge W(M)\), tức \(M^*\) là ghép cặp tốt nhất.

Theo định lý 3, Kuhn năm 1955 và Munkres năm 1957 đưa ra phương pháp sửa
nhãn đỉnh, sao cho ghép cặp cực đại của đồ thị bằng nhau mới dần mở rộng, và
cuối cùng xuất hiện ghép cặp hoàn bị của đồ thị bằng nhau.

\paragraph{Thuật toán Kuhn--Munkres, hay thuật toán KM.}
\begin{enumerate}
  \item Chọn nhãn đỉnh khả thi ban đầu \(l\), và chọn một ghép cặp \(M\)
  trong \(G_l\).
  \item Nếu mọi đỉnh trong \(V_1\) đều liên thuộc với cạnh thuộc \(M\), dừng:
  \(M\) là ghép cặp tốt nhất. Ngược lại, lấy một đỉnh \(u\) của \(G_l\)
  không liên thuộc với cạnh nào trong \(M\), đặt \(S=\{u\}, T=\varnothing\).
  \item Nếu \(P_{G_l}(S)\supset T\), chuyển sang bước 4. Nếu
  \(P_{G_l}(S)=T\), lấy
  \[
    a_l=\min_{x\in S,\ y\notin T}\{l(x)+l(y)-w(xy)\},
  \]
  rồi đặt
  \[
  l'(v)=
  \begin{cases}
  l(v)-a_l, & v\in S,\\
  l(v)+a_l, & v\in T,\\
  l(v), & \text{trường hợp khác}.
  \end{cases}
  \]
  Gán \(l\leftarrow l'\), \(G_l\leftarrow G_{l'}\).
  \item Chọn một đỉnh \(y\in P_{G_l}(S)-T\). Nếu \(y\) liên thuộc với một
  cạnh thuộc \(M\), và \(yz\in M\), đặt
  \(S\leftarrow S\cup\{z\}\), \(T\leftarrow T\cup\{y\}\), rồi quay lại bước
  3. Ngược lại, lấy một đường tăng \(R(u,y)\) của \(M\) trong \(G_l\), đặt
  \(M\leftarrow M\oplus E(R)\), rồi quay lại bước 2.
\end{enumerate}

Dưới đây dùng một ví dụ để hiểu thêm cách thực hiện cụ thể và hiệu quả của
thuật toán. Giả sử hiệu suất của năm nhân viên khi tham gia năm dự án như
sau:
\[
\begin{array}{c|ccccc}
 & Y_1 & Y_2 & Y_3 & Y_4 & Y_5\\ \hline
X_1 & 3 & 5 & 5 & 4 & 1\\
X_2 & 2 & 2 & 0 & 2 & 2\\
X_3 & 2 & 4 & 4 & 1 & 0\\
X_4 & 0 & 1 & 1 & 0 & 0\\
X_5 & 1 & 2 & 1 & 3 & 3
\end{array}
\]
Trước hết, chọn nhãn đỉnh khả thi \(l(v)\) như sau:
\[
\begin{aligned}
l(y_i)&=0,\quad i=1,2,3,4,5,\\
l(x_1)&=\max\{3,5,5,4,1\}=5,\\
l(x_2)&=\max\{2,2,0,2,2\}=2,\\
l(x_3)&=\max\{2,4,4,1,0\}=4,\\
l(x_4)&=\max\{0,1,1,0,0\}=1,\\
l(x_5)&=\max\{1,2,1,3,3\}=3.
\end{aligned}
\]
Lưu ý dòng cuối trong bản in gốc ghi nhầm \(l(x_1)\); theo dữ liệu phải là
\(l(x_5)\).

Xây dựng \(G_l\) và tìm ghép cặp cực đại trên đó. Quá trình như sau:
\[
\begin{array}{c|c|c|l}
\text{Bước} & S & T & M\\ \hline
\text{Chọn ghép cặp đầu} & - & - & \varnothing\\
u=x_1 & \{x_1\} & \varnothing & \varnothing\\
y=y_2 & - & - & \{e(x_1y_2)\}\\
u=x_2 & \{x_2\} & \varnothing & \{e(x_1y_2)\}\\
y=y_1 & - & - & \{e(x_1y_2),e(x_2y_1)\}\\
u=x_3 & \{x_3\} & \varnothing & \{e(x_1y_2),e(x_2y_1)\}\\
y=y_3 & - & - & \{e(x_1y_2),e(x_2y_1),e(x_3y_3)\}\\
u=x_5 & \{x_5\} & \varnothing &
\{e(x_1y_2),e(x_2y_1),e(x_3y_3)\}\\
y=y_5 & - & - &
\{e(x_1y_2),e(x_2y_1),e(x_3y_3),e(x_5y_5)\}\\
u=x_4 & \{x_4\} & \varnothing &
\{e(x_1y_2),e(x_2y_1),e(x_3y_3),e(x_5y_5)\}\\
y=y_3 & \{x_4,x_3\} & \{y_3\} &
\{e(x_1y_2),e(x_2y_1),e(x_3y_3),e(x_5y_5)\}\\
y=y_2 & \{x_4,x_3,x_1\} & \{y_3,y_2\} &
\{e(x_1y_2),e(x_2y_1),e(x_3y_3),e(x_5y_5)\}\\
P_{G_l}(S)=T & - & - & \text{không có ghép cặp hoàn bị}
\end{array}
\]

Ghép cặp cực đại cuối cùng có thể biểu diễn bằng các cạnh
\[
  x_1-y_2,\qquad x_2-y_1,\qquad x_3-y_3,\qquad x_5-y_5.
\]
Các cạnh này là các cạnh được tô đậm trong hình của bản gốc. Trong \(G_l\)
không có ghép cặp hoàn bị, vì vậy ta sửa nhãn đỉnh.

Do \(u=x_4\), \(S=\{x_4,x_3,x_1\}\), \(T=\{y_3,y_2\}\), nên
\[
  a_l=\min_{x\in S,\ y\notin T}\{l(x)+l(y)-w(xy)\}=1.
\]
Nhãn đỉnh sau khi sửa là:
\[
\begin{gathered}
x_1=4,\quad x_2=2,\quad x_3=3,\quad x_4=0,\quad x_5=3,\\
y_1=0,\quad y_2=1,\quad y_3=1,\quad y_4=0,\quad y_5=0.
\end{gathered}
\]
Theo nhãn đỉnh mới, xây dựng \(G_l\) và tìm được một ghép cặp hoàn bị:
\[
  x_1-y_4,\qquad x_2-y_1,\qquad x_3-y_3,\qquad x_4-y_2,\qquad x_5-y_5.
\]
Đây là ghép cặp tốt nhất; trọng số lớn nhất của nó là
\[
  4+2+4+1+3=14.
\]
Bài toán được giải xong.

Qua quá trình giải và so sánh mô tả thuật toán KM với thuật toán Hungary, có
thể thấy ngoài phần sửa nhãn đỉnh ở bước thứ ba của KM, phần lớn hai thuật
toán là giống nhau. Có thể nói thuật toán KM chính là thuật toán Hungary
được bổ sung bước sửa nhãn đỉnh. Vì \(M\) được mở rộng dần, thông tin sinh
ra trong mỗi quá trình tăng trên \(G_l\) đều được kế thừa và sử dụng hợp lý.
Do đó, mặc dù thuật toán KM đưa trọng số vào xét và có phạm vi áp dụng rộng
hơn, về bậc độ phức tạp thời gian, nó về bản chất giống thuật toán Hungary,
cũng là \(O(|E(G)|)\) cho mỗi quá trình mở rộng trong cách trình bày này.

\section{Kết luận}

Đồ thị hai phía là một loại đồ thị đặc biệt. Với một đồ thị hai phía \(G\),
tập đỉnh \(V(G)\) của nó có thể chia thành hai tập con bù nhau \(V_1,V_2\);
với hai đỉnh \(x,y\) thuộc cùng một tập con, luôn có
\[
  e(x,y)\notin E(G).
\]
Nhờ tính đặc biệt này, đồ thị hai phía được ứng dụng rộng rãi như một mô hình
hữu hiệu để mô tả quan hệ giữa hai loại sự vật khác nhau trong thế giới thực.

Cấu trúc đồ thị thường phản ánh thuộc tính bản chất của bài toán nhờ lượng
thông tin đa dạng mà nó chứa, từ đó nắm được điểm then chốt của bài toán. Đặc
biệt khi giải quyết các vấn đề như tính hậu hiệu và thảm họa số chiều trong
quy hoạch động, mô hình đồ thị thường phát huy hiệu quả tốt. Lại vì các thuật
toán xây dựng trên đồ thị hai phía có nhiều ưu thế so với thuật toán trên đồ
thị tổng quát, việc chuyển một đồ thị tổng quát thành dạng đồ thị hai phía có
tác dụng rất lớn trong lập trình. Mấu chốt của việc chuyển đổi thường là
xuất phát từ chính điều kiện của đề bài, khai thác thông tin sâu hơn trong
đề, rồi phân loại các đỉnh của đồ thị có quan hệ phức tạp. Chẳng hạn trên
bàn cờ, tô màu thường là một cách phân loại rất tốt.

Sau khi xây dựng được mô hình đồ thị hai phía, việc giải cũng rất quan
trọng. Bài viết này đã đưa ra một số thuật toán cơ bản và định lý liên quan
của đồ thị hai phía. Sử dụng hợp lý các thông tin ấy về cơ bản có thể giải
quyết phần lớn bài toán đồ thị hai phía gặp trong lập trình. Nếu có thêm
nhiều hạn chế và yêu cầu hơn, ta có thể dùng các mô hình phức tạp hơn như
luồng cực đại để giải. Nhưng xét về hiệu quả thuật toán và độ phức tạp lập
trình, thuật toán dựa trên đồ thị hai phía thường đơn giản và hiệu quả hơn
thuật toán dựa trên luồng cực đại.

Đối với thuật toán ghép cặp, ngoài ghép cặp của đồ thị hai phía, ghép cặp
trên đồ thị tổng quát cũng có định nghĩa, nhưng thuật toán giải tương đối
phức tạp. Đồng thời, theo quan điểm cá nhân của tôi, vì ghép cặp là quan hệ
một-một giữa hai vật, nên việc xây dựng quan hệ ghép cặp giữa hai loại sự
vật, tức ghép cặp trên đồ thị hai phía, có giá trị thực dụng cao hơn ghép cặp
trên đồ thị tổng quát.

Tóm lại, tôi cho rằng đồ thị hai phía là một mô hình hiệu quả và có giá trị
sử dụng rộng rãi. Sử dụng mô hình đồ thị hai phía một cách hợp lý và hiệu quả
sẽ nâng cao đáng kể năng lực lập trình cũng như năng lực giải quyết các vấn
đề thực tế trong đời sống.

\appendix
\section{Chứng minh định lý hôn nhân Hall}

Chứng minh định lý hôn nhân Hall như sau.

Rõ ràng, với mọi tập con \(A\) của \(X\), điều kiện \(|A|\le |P(A)|\) là
điều kiện cần để tồn tại ghép cặp hoàn bị từ \(X\) vào \(Y\). Vì nếu có một
tập con \(A_0\) của \(X\) sao cho \(|A_0|>|P(A_0)|\), thì không thể xác định
một quan hệ một-một giữa các đỉnh của \(A_0\) và các đỉnh trong một tập con
của \(Y\). Sau đây dùng hai cách để chứng minh điều kiện này là đủ.

\subsection{Chứng minh bằng quy nạp}

Với mọi đồ thị hai phía mà \(X\) chỉ có một đỉnh, điều kiện hiển nhiên là
đủ, vì đỉnh ấy nhất định kề với ít nhất một đỉnh trong \(Y\). Bây giờ giả sử
đối với mọi đồ thị hai phía có không quá \(m-1\) đỉnh trong \(X\), điều kiện
này là đủ. Xét một đồ thị hai phía \(G\), trong đó \(X\) có \(m\) đỉnh và
mọi tập con \(A\) của \(X\) đều thỏa \(|A|\le |P(A)|\). Ta xét hai trường
hợp sau để chứng minh tồn tại ghép cặp hoàn bị từ \(X\) vào \(Y\).

\paragraph{Trường hợp 1.}
Với mọi tập con thực sự không rỗng \(A\) của \(X\), đều có
\(|A|<|P(A)|\). Ta có thể lấy một đỉnh bất kỳ \(x_0\in X\) và ghép nó với
một đỉnh \(y_0\in Y\). Xóa khỏi \(G\) hai đỉnh \(x_0,y_0\) cùng mọi cạnh
liên thuộc với chúng, thu được đồ thị \(G'\). Trong đồ thị hai phía \(G'\),
mỗi tập con \(A\) của \(X-\{x_0\}\) vẫn kề với ít nhất \(|A|\) đỉnh trong
\(Y-\{y_0\}\). Do đó, theo giả thiết quy nạp, tồn tại một ghép cặp hoàn bị
từ \(X-\{x_0\}\) vào \(Y-\{y_0\}\). Vì vậy trong đồ thị \(G\) tồn tại ghép
cặp hoàn bị từ \(X\) vào \(Y\).

\paragraph{Trường hợp 2.}
Tồn tại một tập con thực sự không rỗng \(A_0\) của \(X\) sao cho
\(|A_0|=|P(A_0)|\). Gọi \(G'\) là đồ thị con của \(G\) chứa tập đỉnh
\(A_0\), tập đỉnh \(P(A_0)\) kề với các đỉnh trong \(A_0\), và toàn bộ các
cạnh nối hai tập đỉnh này. Vì mọi tập con \(A\) của \(A_0\) kề với ít nhất
\(|A|\) đỉnh trong \(P(A_0)\), theo giả thiết quy nạp, trong \(G'\) tồn tại
một ghép cặp hoàn bị từ \(A_0\) vào \(P(A_0)\).

Gọi \(G''\) là đồ thị con chứa các tập đỉnh \(X-A_0\) và \(Y-P(A_0)\), cùng
các cạnh nối hai tập này. Ta khẳng định trong \(G''\) có ghép cặp hoàn bị từ
\(X-A_0\) vào \(Y-P(A_0)\). Nếu không, theo giả thiết quy nạp, sẽ tồn tại một
tập con \(A_1\) của \(X-A_0\) sao cho trong \(G''\), các đỉnh của \(A_1\) kề
với ít hơn \(|A_1|\) đỉnh trong \(Y-P(A_0)\). Điều đó có nghĩa trong đồ thị
\(G\), các đỉnh của tập con \(A_0\cup A_1\) kề với ít hơn
\(|A_0\cup A_1|\) đỉnh, mâu thuẫn với giả thiết ``với mọi tập con \(A\) của
\(X\), đều có \(|A|\le |P(A)|\)''. Vì thế trong \(G''\) tồn tại ghép cặp
hoàn bị từ \(X-A_0\) vào \(Y-P(A_0)\). Từ đây suy ra trong đồ thị \(G\) có
ghép cặp hoàn bị từ \(X\) vào \(Y\). Chứng minh hoàn tất.

\subsection{Chứng minh bằng phản chứng}

Giả sử \(G\) không tồn tại ghép cặp hoàn bị. Khi đó \(n=|V_1|\) đỉnh của
\(V_1\) không thể ghép đôi với \(n\) đỉnh khác nhau của \(V_2\). Bây giờ
giả sử đã tìm được một ghép cặp có số cạnh lớn nhất \(M^*\) của \(G\). Khi
đó trong \(V_1\) ít nhất có một đỉnh không được \(M^*\) bão hòa, nếu không
\(M^*\) đã là ghép cặp hoàn bị. Gọi đỉnh đó là \(v_0\). Từ \(v_0\), ta tìm
tất cả các đường luân phiên đối với \(M^*\), và gọi tập đỉnh trên các đường
luân phiên ấy là \(S\). Ký hiệu tập đỉnh thuộc \(V_1\) trong \(S\) là \(A\),
tập đỉnh thuộc \(V_2\) trong \(S\) là \(B\), tức
\[
  S=A\cup B,\qquad A=S\cap V_1,\qquad B=S\cap V_2.
\]

Theo tính chất của đường luân phiên, mọi đỉnh kề với đỉnh trong \(A\) đều
phải là đỉnh được \(M^*\) bão hòa; các đỉnh ấy chắc chắn nằm trên một đường
luân phiên đối với \(M^*\) bắt đầu từ \(v_0\). Vì vậy
\[
  P(A)=B. \tag{1}
\]
Ngoài ra, trong \(A\), trừ \(v_0\), mỗi đỉnh đều ghép đôi một-một với một
đỉnh trong \(B\), nên
\[
  |B|=|A|-1. \tag{2}
\]
Từ (1) và (2) suy ra
\[
  |P(A)|=|A|-1<|A|,
\]
mâu thuẫn với điều kiện giả thiết. Định lý được chứng minh.

\section{Tài liệu tham khảo}

\begin{enumerate}
  \item \textit{Tổ hợp học và thuật toán}, Nhà xuất bản Đại học Khoa học và
  Công nghệ Trung Quốc, Yang Zhensheng biên soạn.
  \item \textit{Lý thuyết đồ thị và thuật toán}, Nhà xuất bản Công nghiệp
  Hàng không, Xiao Weishu chủ biên.
  \item Tạp chí quý \textit{Olympic Tin học}, do Ủy ban tổ chức vòng liên
  tỉnh Olympic Tin học thiếu niên toàn quốc chủ trì.
  \item \textit{Giáo trình thi mô hình hóa toán học}, Nhà xuất bản Giáo dục
  Giang Tô, Li Shangzhi chủ biên.
\end{enumerate}

Ngoài ra, bài viết còn tham khảo các đề và lời giải liên quan của BOI, USACO
và ACM, đồng thời tổng hợp kết quả thảo luận với các học sinh khác ở tỉnh
Giang Tô.

\section{Phần trình chiếu kèm trong PDF}

Sau bài viết chính, bản PDF còn chứa một bản trình chiếu ngắn. Nội dung dưới
đây chuyển toàn bộ phần chữ trích được từ các trang đó sang tiếng Việt.

\subsection*{Trang bìa}

\begin{center}
\Large Thuật toán và ứng dụng của đồ thị hai phía

\normalsize Trường Trung học trực thuộc Đại học Sư phạm Nam Kinh

Sun Fangcheng
\end{center}

\subsection*{Mục lục}

\begin{enumerate}
  \item Khái niệm ghép cặp.
  \item Định nghĩa và nhận biết đồ thị hai phía.
  \item Ghép cặp cực đại của đồ thị hai phía.
  \item Bài toán tập phủ nhỏ nhất của đồ thị hai phía.
  \item Bài toán ghép cặp tốt nhất của đồ thị hai phía.
  \item Kết luận.
\end{enumerate}

\subsection*{Khái niệm ghép cặp}

Cho đồ thị \(G=(V(G),E(G))\), và \(M\) là một tập con của \(E(G)\). Nếu hai
cạnh bất kỳ trong \(M\) đều không kề nhau, thì \(M\) được gọi là một ghép
cặp của \(G\). Hai đầu mút của một cạnh trong \(M\) được gọi là được ghép đôi
dưới \(M\).

Nếu một cạnh nào đó của ghép cặp \(M\) liên thuộc với đỉnh \(v\), thì \(M\)
bão hòa đỉnh \(v\), và \(v\) là đỉnh bão hòa đối với \(M\).

Giả sử \(M\) là một ghép cặp của đồ thị \(G\). Nếu trong \(G\) tồn tại một
đường đi đơn \(R\) mà các cạnh trên đường đi luân phiên là cạnh ghép thuộc
\(M\) và cạnh không ghép không thuộc \(M\), thì \(R\) là đường luân phiên.
Nếu hai đầu mút của \(R\) đều là đỉnh chưa bão hòa đối với \(M\), thì \(R\)
là đường tăng.

Cho đồ thị \(G=(V(G),E(G))\), trong đó \(V(G)\) được chia thành hai tập đỉnh
con không rỗng và bù nhau \(X\) và \(Y\). Nếu một ghép cặp
\(M\subseteq E(G)\) của \(G\) bão hòa mọi đỉnh trong \(X\), thì \(M\) là
ghép cặp hoàn bị của \(G\).

\subsection*{Định nghĩa và nhận biết đồ thị hai phía}

Cho đồ thị \(G=(V,E)\). Nếu có thể chia \(V\) thành hai tập \(V_1,V_2\), sao
cho mỗi cạnh trong \(E\) có một đầu mút ở \(V_1\) và đầu mút kia ở \(V_2\),
thì đồ thị ấy là đồ thị hai phía. Đồ thị hai phía cũng có thể viết là
\(G=(V_1,V_2;E)\). Với \(V\subseteq V_1\), dùng \(P(V)\) để chỉ tập các đỉnh
trong \(V_2\) kề với \(V\).

Nếu trong đồ thị hai phía \(G\), mỗi đỉnh của \(V_1\) đều kề với mọi đỉnh
của \(V_2\), thì \(G\) là đồ thị hai phía đầy đủ.

Định lý nhận biết: một đồ thị vô hướng \(G\) là đồ thị hai phía khi và chỉ
khi mọi chu trình của \(G\) đều có độ dài chẵn. Nếu không có chu trình, xem
độ dài chu trình là \(0\), và \(0\) là số chẵn.

\subsection*{Ghép cặp cực đại của đồ thị hai phía}

Edmonds năm 1965 đưa ra thuật toán Hungary để giải ghép cặp cực đại của đồ
thị hai phía:
\begin{enumerate}
  \item Lấy một ghép cặp ban đầu \(M\) trong \(G\).
  \item Nếu mọi đỉnh trong \(X\) đều là đầu mút của cạnh thuộc \(M\), dừng:
  \(M\) là ghép cặp hoàn bị. Ngược lại, lấy một đỉnh \(u\) trong \(X\) không
  liên thuộc với cạnh nào trong \(M\), đặt \(S=\{u\},T=\varnothing\).
  \item Nếu \(P(S)=T\), dừng: không có ghép cặp hoàn bị. Ngược lại, lấy
  \(y\in P(S)-T\).
  \item Nếu \(y\) là đầu mút của cạnh thuộc \(M\), giả sử \(yz\in M\), đặt
  \(S\leftarrow S\cup\{z\}\), \(T\leftarrow T\cup\{y\}\), rồi quay lại bước
  3. Ngược lại, lấy đường tăng \(R(u,y)\), đặt
  \(M\leftarrow M\oplus E(R)\), rồi quay lại bước 2.
\end{enumerate}

\subsection*{Tập phủ nhỏ nhất và ghép cặp cực đại}

Bài toán tập phủ nhỏ nhất của đồ thị tổng quát là một bài toán NPC đã được
chứng minh. Nói cách khác, bài toán tập phủ nhỏ nhất của đồ thị tổng quát
không có thuật toán hiệu quả đã biết. Vì vậy, nếu một bài toán thực tế chỉ
được trừu tượng hóa thành tập phủ nhỏ nhất của đồ thị tổng quát, mô hình đó
khó đem lại giá trị thuật toán.

Nhưng nếu bài toán có thể được trừu tượng hóa thành tập phủ nhỏ nhất của đồ
thị hai phía, kết cục sẽ khác. Do tính đặc thù của đồ thị hai phía, bài toán
tập phủ nhỏ nhất của đồ thị hai phía có thuật toán đa thức.

Trong chứng minh quan hệ này cần dùng định lý hôn nhân Hall: cho đồ thị hai
phía \(G\), tập đỉnh chia thành \(V_1\) và \(V_2\). Trong \(G\) tồn tại ghép
cặp hoàn bị đối với \(V_1\) khi và chỉ khi với mọi \(S\subseteq V_1\), đều
có \(|S|\le |P(S)|\).

Năm 1931, König đưa ra định lý về quan hệ giữa ghép cặp cực đại và phủ nhỏ
nhất: trong đồ thị hai phía \(G\), nếu \(M^*\) là ghép cặp cực đại và \(K^*\)
là tập phủ nhỏ nhất, thì
\[
  |M^*|=|K^*|.
\]

\subsection*{Bài toán ghép cặp tốt nhất}

Cho đồ thị hai phía đầy đủ có trọng số
\[
G=(V_1,V_2;E),\qquad
V_1=\{x_1,x_2,\ldots,x_n\},\qquad
V_2=\{y_1,y_2,\ldots,y_n\},
\]
với trọng số \(w(x_iy_j)\ge 0\). Nếu một ghép cặp hoàn bị \(M\) thỏa
\(W(M)\ge W(M')\) với mọi ghép cặp hoàn bị \(M'\), thì \(M\) là ghép cặp tốt
nhất của \(G\).

Vì có trọng số, ghép cặp tốt nhất không thể trực tiếp dùng thuật toán ghép
cặp cực đại để giải. Đồng thời, định nghĩa ghép cặp tốt nhất được xây dựng
trên đồ thị hai phía đầy đủ có trọng số; với đồ thị không đầy đủ, có thể
thêm các cạnh trọng số \(0\), hoặc giá trị khác không ảnh hưởng đến kết quả,
để biến đồ thị hai phía thành đồ thị hai phía đầy đủ rồi dùng thuật toán
ghép cặp tốt nhất.

Trước khi giới thiệu thuật toán KM, cần nêu khái niệm sau. Ánh xạ
\(l:V(G)\to R\) thỏa
\[
  \forall x\in V_1,\quad \forall y\in V_2,\quad
  l(x)+l(y)\ge w(xy)
\]
được gọi là nhãn đỉnh khả thi của đồ thị hai phía \(G\). Đặt
\[
  E_l=\{xy\mid xy\in E(G),\ l(x)+l(y)=w(xy)\}.
\]
Đồ thị con sinh lấy \(E_l\) làm tập cạnh được gọi là đồ thị bằng nhau, ký
hiệu \(G_l\). Có thể chứng minh rằng ghép cặp hoàn bị của \(G_l\) chính là
ghép cặp tốt nhất của \(G\).

Trên cơ sở đó, Kuhn năm 1955 và Munkres năm 1957 đưa ra phương pháp sửa nhãn
đỉnh, khiến ghép cặp cực đại của đồ thị bằng nhau mới dần mở rộng, cuối cùng
xuất hiện ghép cặp hoàn bị của đồ thị bằng nhau. Đó là thuật toán KM.

\subsection*{Thuật toán KM}

\begin{enumerate}
  \item Chọn nhãn đỉnh khả thi ban đầu \(l\), và chọn một ghép cặp ban đầu
  \(M\) trên \(G_l\).
  \item Nếu mọi đỉnh trong \(V_1\) đều là đầu mút của cạnh thuộc \(M\), dừng:
  \(M\) là ghép cặp tốt nhất. Ngược lại, lấy một đỉnh \(u\) trong \(G_l\)
  không liên thuộc với cạnh nào trong \(M\), đặt \(S=\{u\},T=\varnothing\).
  \item Nếu \(P(S)\supset T\), chuyển sang bước 4. Nếu \(P(S)=T\), lấy
  \[
  a_l=\min_{x\in S,\ y\notin T}\{l(x)+l(y)-w(xy)\},
  \]
  rồi sửa nhãn theo
  \[
  l'(v)=
  \begin{cases}
  l(v)-a_l, & v\in S,\\
  l(v)+a_l, & v\in T,\\
  l(v), & \text{trường hợp khác}.
  \end{cases}
  \]
  Sau đó đặt \(l\leftarrow l'\), \(G_l\leftarrow G_{l'}\).
  \item Chọn một đỉnh \(y\in P(S)-T\). Nếu \(y\) là đầu mút của cạnh thuộc
  \(M\), và \(yz\in M\), đặt \(S\leftarrow S\cup\{z\}\),
  \(T\leftarrow T\cup\{y\}\), rồi quay lại bước 3. Ngược lại, lấy đường tăng
  \(R(u,y)\) của \(M\) trong \(G_l\), đặt \(M\leftarrow M\oplus E(R)\), rồi
  quay lại bước 2.
\end{enumerate}

\subsection*{Ví dụ trên trang trình chiếu}

Một công ty có các nhân viên \(x_1,x_2,\ldots,x_n\), họ cần làm các công
việc \(y_1,y_2,\ldots,y_n\). Mỗi người có thể làm một vài công việc và có
hiệu suất cố định đối với từng công việc. Cần tìm phương án phân công sao
cho tổng hiệu suất cao nhất; một người chỉ làm một công việc, và một công
việc phải do một người độc lập hoàn thành. Không yêu cầu mọi người đều có
việc.

Ví dụ dữ liệu:
\[
\begin{array}{c|ccccc}
 & Y_1 & Y_2 & Y_3 & Y_4 & Y_5\\ \hline
X_1 & 3 & 5 & 5 & 4 & 1\\
X_2 & 2 & 2 & 0 & 2 & 2\\
X_3 & 2 & 4 & 4 & 1 & 0\\
X_4 & 0 & 1 & 1 & 0 & 0\\
X_5 & 1 & 2 & 1 & 3 & 3
\end{array}
\]
Nếu công nhân \(x\) hoàn toàn không thể tham gia công việc \(y\), đặt
\(w(x,y)=0\).

Chọn nhãn đỉnh khả thi:
\[
\begin{aligned}
l(y_i)&=0,\quad i=1,2,3,4,5,\\
l(x_1)&=\max\{3,5,5,4,1\}=5,\\
l(x_2)&=\max\{2,2,0,2,2\}=2,\\
l(x_3)&=\max\{2,4,4,1,0\}=4,\\
l(x_4)&=\max\{0,1,1,0,0\}=1,\\
l(x_5)&=\max\{1,2,1,3,3\}=3.
\end{aligned}
\]
Xây dựng \(G_l\), tìm ghép cặp cực đại; quá trình chi tiết đã nằm trong bài
viết chính.

Ghép cặp cực đại ban đầu có thể biểu diễn bởi
\[
  x_1-y_2,\qquad x_2-y_1,\qquad x_3-y_3,\qquad x_5-y_5.
\]
Trong \(G_l\) chưa có ghép cặp hoàn bị, nên sửa nhãn đỉnh. Do
\[
  u=x_4,\qquad S=\{x_4,x_3,x_1\},\qquad T=\{y_3,y_2\},
\]
nên
\[
  a_l=\min_{x\in S,\ y\notin T}\{l(x)+l(y)-w(xy)\}=1.
\]
Nhãn đỉnh sau khi sửa:
\[
\begin{gathered}
x_1=4,\quad x_2=2,\quad x_3=3,\quad x_4=0,\quad x_5=3,\\
y_1=0,\quad y_2=1,\quad y_3=1,\quad y_4=0,\quad y_5=0.
\end{gathered}
\]
Theo nhãn mới, \(G_l\) có một ghép cặp hoàn bị, cho ghép cặp tốt nhất với
trọng số lớn nhất
\[
  4+2+4+1+3=14.
\]

\subsection*{Kết luận của trang trình chiếu}

Đồ thị hai phía là một loại đồ thị đặc biệt. Nó không chỉ có ưu điểm chung
của mô hình đồ thị là chứa lượng thông tin phong phú, mà còn có nhiều nội hàm
và ưu thế thuật toán mà đồ thị tổng quát không có.

Các đỉnh của đồ thị hai phía được chia thành hai phần; điều này có liên hệ
sâu sắc với quan hệ tương ứng, hay quan hệ ghép cặp, trong tự nhiên và trong
toán học. Vì vậy, thuật toán ghép cặp là lõi của mọi thuật toán trên đồ thị
hai phía.

Nếu có thể trừu tượng hóa một bài toán thực tế thành mâu thuẫn giữa hai loại
sự vật, thì bài toán ấy có thể được mô hình hóa bằng đồ thị hai phía. Do đó
mô hình đồ thị hai phía có ứng dụng rộng rãi. Đồng thời, thuật toán trên đồ
thị hai phía có đặc điểm hiệu quả và thực dụng, khiến việc giải bài bằng mô
hình đồ thị hai phía trở nên khả thi.

Tóm lại, tôi cho rằng đồ thị hai phía là một mô hình hiệu quả và có giá trị
sử dụng rộng rãi. Sử dụng mô hình đồ thị hai phía một cách hợp lý và hiệu quả
sẽ nâng cao đáng kể năng lực lập trình cũng như năng lực giải quyết các vấn
đề thực tế trong đời sống.

\begin{center}
Cảm ơn!
\end{center}

\end{document}
